\chapter{Animations and easing}
\label{ch:animations}

\begin{aside}
Snap-changes feel sudden; tweened changes feel polished. A
quarter-second fade is the difference between a popup that
appears and a popup that arrives.
\end{aside}

\section{The interpolator}

An animation maps time to a value:

\begin{lstlisting}
auto fade = animate(300ms, [](float t) {
    return lerp(0.0f, 1.0f, t);
});
\end{lstlisting}

The runtime ticks $t$ from 0 to 1 over 300 ms and calls the
lambda. The return value drives a signal which drives a
re-render.

\section{Easing}

Linear time often feels mechanical. Easing curves shape the
$t$ value:

\begin{itemize}[leftmargin=*]
  \item \code{ease\_in}: slow start, fast finish.
  \item \code{ease\_out}: fast start, slow finish.
  \item \code{ease\_in\_out}: slow at both ends.
  \item \code{spring}: overshoot and settle.
  \item \code{bounce}: ballistic return.
\end{itemize}

The cubic-bezier formulation:

\begin{lstlisting}
float ease_in_out(float t) {
    return t < 0.5f
         ? 2 * t * t
         : 1 - std::pow(-2*t + 2, 2) / 2;
}
\end{lstlisting}

\section{Composing}

Animations compose:

\begin{itemize}[leftmargin=*]
  \item Sequential: A then B then C.
  \item Parallel: A and B simultaneously.
  \item Conditional: A if X, else B.
\end{itemize}

\maya{}'s \code{seq()}, \code{par()}, \code{when()}
combinators build these from primitives.

\section{Frame rate}

The runtime ticks at a fixed rate (60 Hz default). Each
animation gets a frame; multiple animations share one
frame.

A long-running animation does not block the loop --- it
schedules a wake at its next tick deadline and otherwise
sleeps.

\section{Interrupting}

An animation that is interrupted (state changes mid-tween)
either:

\begin{itemize}[leftmargin=*]
  \item Snaps to current state and starts a new tween from
    there.
  \item Reverses (good for ``open / close'' animations).
  \item Cancels.
\end{itemize}

The choice is per-animation. Snapping causes visual jolts;
reversing is smoother; cancelling leaves the value in an
intermediate state.

\section{Performance}

A frame with 100 active animations runs 100 callbacks per
tick. Each callback is cheap (one arithmetic op typically)
but they add up. For very large animation counts, batch:
one callback that updates many values.

\section{Disabled motion}

Some users (vestibular disorders, terminal in motion on a
phone) prefer reduced motion. Respect:

\begin{itemize}[leftmargin=*]
  \item \code{REDUCE\_MOTION=1} environment variable.
  \item Application setting.
\end{itemize}

When set, animations snap to end state instantly.

\begin{warning}
Animations are a polish, not a function. Snap-mode must
always be a viable path; never make UI rely on the
in-between frames for information.
\end{warning}

\section{Terminal-side animations}

Some image protocols (Kitty) support terminal-side
animations: emit frames and timing, terminal cycles. CPU on
the application side is zero.

For a spinner or a progress beachball, terminal-side is
ideal --- the app sleeps and the screen still moves.

\section{Recap}

\begin{recap}
\begin{itemize}[leftmargin=*]
  \item Animation = time-to-value interpolator.
  \item Easing curves shape perceived feel.
  \item Compose with seq, par, when.
  \item Interruption strategy is per-animation.
  \item Respect reduced-motion preferences.
\end{itemize}
\end{recap}

\section*{Workshop}

\begin{enumerate}[leftmargin=*]
  \item Animate a panel slide-in. Try ease-in-out and
    spring; pick what feels right.
  \item Add reduced-motion support. Test with the env var.
  \item Build a sequential animation chain (fade in,
    delay, fade out).
\end{enumerate}
